\documentclass[11pt,a4paper]{article}
\usepackage[margin=0.95in]{geometry}
\usepackage[T1]{fontenc}
\usepackage{lmodern}
\usepackage{amsmath,amssymb}
\usepackage{booktabs}
\usepackage{xcolor}
\usepackage{fancyhdr}
\usepackage{hyperref}
\setlength{\parskip}{0.46em}
\setlength{\parindent}{0pt}
\setlength{\headheight}{15pt}
\linespread{1.07}
\pagestyle{fancy}
\fancyhf{}
\lhead{CST444 Soft Computing - Module 5}
\rhead{Exam-Ready Beginner Guide}
\cfoot{\thepage}
\hypersetup{colorlinks=true,linkcolor=blue,urlcolor=blue}

\newcommand{\examtip}[1]{%
\vspace{0.2em}
\noindent\fcolorbox{black}{gray!12}{\parbox{0.97\linewidth}{\textbf{Exam Tip:} #1}}%
\vspace{0.2em}
}

\begin{document}

\section*{Module 5 Topic Map from OCR Questions}
\textbf{Frequently repeated concepts}
\begin{itemize}
\item Multi-objective optimization problem (MOOP): formulation and examples.
\item Linear vs nonlinear MOOP; convex vs non-convex MOOP.
\item Dominance relation and its properties.
\item Pareto optimality, non-dominated set, Pareto front.
\item Procedure for finding non-dominated set.
\item Neuro-fuzzy hybrid systems: classification, characteristics, working.
\item Genetic-neuro hybrid systems: block diagram, steps, advantages.
\end{itemize}

\textbf{High-mark areas}
\begin{itemize}
\item Convex/non-convex MOOP + non-dominated set derivation.
\item Properties of dominance relation and Pareto optimality.
\item Genetic-neuro and neuro-fuzzy hybrid architecture + advantages.
\end{itemize}

\textbf{Must-study-first order}
\begin{enumerate}
\item MOOP formulation and objective/decision spaces.
\item Dominance and Pareto front.
\item Convexity concepts.
\item Non-dominated sorting intuition.
\item Genetic-neuro hybrid flow.
\item Neuro-fuzzy hybrid classification and features.
\end{enumerate}

\section*{Core Theory (Beginner-Friendly, Exam-Formal)}

\subsection*{1) Multi-Objective Optimization Problem (MOOP)}
Real-world design often needs simultaneous optimization of conflicting objectives.

\textbf{General formulation}
\[
\min/\max\ f_i(x),\ i=1,2,\dots,m
\]
subject to constraints
\[
g_j(x)\le 0,\quad j=1,2,\dots,J,
\qquad
h_k(x)=0,\quad k=1,2,\dots,K
\]
with decision vector $x=[x_1,x_2,\dots,x_n]$.

\textbf{Key interpretation}
\begin{itemize}
\item Decision space contains candidate vectors $x$.
\item Objective space contains mapped vectors $f(x)$.
\item MOOP typically yields a set of trade-off optimal solutions, not a single point.
\end{itemize}

\subsection*{2) Linear/Nonlinear and Convex/Non-convex MOOP}
\textbf{Linear MOOP:} all objectives and constraints are linear.

\textbf{Nonlinear MOOP:} at least one objective or constraint is nonlinear.

\textbf{Convex MOOP:} convex objectives and convex feasible region.
\[
f(\lambda x+(1-\lambda)y)\le \lambda f(x)+(1-\lambda)f(y),\quad 0<\lambda<1
\]

\textbf{Non-convex MOOP:} convexity condition does not hold globally.

\subsection*{3) Dominance Relation}
Solution $x^1$ dominates $x^2$ if:
\begin{itemize}
\item $x^1$ is no worse than $x^2$ in all objectives.
\item $x^1$ is strictly better in at least one objective.
\end{itemize}

\textbf{Properties (as in notes)}
\begin{itemize}
\item Not reflexive.
\item Not symmetric.
\item Not antisymmetric (in the strict dominance sense used here).
\item Transitive.
\end{itemize}

\subsection*{4) Pareto Optimality and Non-dominated Set}
\begin{itemize}
\item Non-dominated set in a set $P$: all members not dominated by any other member in $P$.
\item Pareto-optimal set: non-dominated set of the entire feasible search space.
\item Pareto front: mapped boundary of Pareto-optimal set in objective space.
\end{itemize}

\begin{center}
\framebox[0.97\linewidth][c]{
\begin{minipage}{0.93\linewidth}
\vspace{0.08cm}
\centerline{\textbf{Pareto Front Idea}}
\centerline{Improve one objective $\Rightarrow$ at least one other objective worsens on the front.}
\vspace{0.08cm}
\end{minipage}}
\end{center}

\subsection*{5) Hybrid Systems in Module 5}
A hybrid system combines at least two intelligent paradigms (fuzzy logic, neural network, genetic algorithm, etc.) to gain complementary strengths.

\textbf{Types listed}
\begin{itemize}
\item Neuro-fuzzy hybrid
\item Neuro-genetic (genetic-neuro) hybrid
\item Fuzzy-genetic hybrid
\end{itemize}

\subsection*{6) Genetic-Neuro Hybrid System}
GA optimizes neural network components such as weights, topology, and control parameters.

\begin{center}
\framebox[0.97\linewidth][c]{
\begin{minipage}{0.93\linewidth}
\vspace{0.08cm}
\centerline{\textbf{Genetic-Neuro Block Flow}}
\centerline{Initial Population $\rightarrow$ Selection/Crossover/Mutation $\rightarrow$ ANN Evaluation}
\centerline{$\rightarrow$ Fitness Calculation $\rightarrow$ Parameter Update $\rightarrow$ Repeat}
\vspace{0.08cm}
\end{minipage}}
\end{center}

\textbf{Advantages from notes}
\begin{itemize}
\item Topology optimization for ANN.
\item Weight/parameter optimization.
\item Can optimize learning rate, momentum, tolerance, etc.
\item Mimics human decision style in complex tasks.
\end{itemize}

\subsection*{7) Neuro-Fuzzy Hybrid System}
Neuro-fuzzy is a fuzzy system trained using neural learning principles.

\begin{center}
\framebox[0.97\linewidth][c]{
\begin{minipage}{0.93\linewidth}
\vspace{0.08cm}
\centerline{\textbf{Neuro-Fuzzy Structure}}
\centerline{Crisp Input $\rightarrow$ Fuzzification $\rightarrow$ Rule Layer/Neural Learning $\rightarrow$ Defuzzification $\rightarrow$ Crisp Output}
\vspace{0.08cm}
\end{minipage}}
\end{center}

\textbf{Characteristics from notes}
\begin{itemize}
\item Handles numeric and linguistic information.
\item Works with uncertain/imprecise information.
\item Self-learning, self-organizing, self-tuning behavior.
\item Supports conflict resolution by aggregation/collaboration.
\end{itemize}

\examtip{In long answers, always separate ``definition'', ``properties'', ``diagram'', and ``application'' as four clear chunks to maximize marks quickly.}

\section*{Numericals and Derivations (Stepwise)}

\subsection*{N1) Dominance check using air-ticket style example}
\textbf{Given}
Objectives are both minimization: (Time, Cost).
Sample points: A(2,7.5), B(3,6), C(3,7.5), D(4,5), F(5,4.5).

\textbf{Find}
Who dominates C? Build one non-dominated set.

\textbf{Rule}
$x$ dominates $y$ if $x$ is no worse in all objectives and strictly better in at least one.

\textbf{Checks}
\begin{itemize}
\item B vs C: Time equal (3=3), cost better (6<7.5) $\Rightarrow$ B dominates C.
\item A vs C: Time better (2<3), cost equal (7.5=7.5) $\Rightarrow$ A dominates C.
\end{itemize}

\textbf{Non-dominated set example}
\[
\{A,B,D,F\}
\]
(no pair in this set dominates another in both objectives).

\subsection*{N2) Convexity condition check (symbolic)}
\textbf{Given}
A function candidate $f$ and two points $x,y$.

\textbf{Find}
Convexity test condition.

\textbf{Formula}
\[
f(\lambda x+(1-\lambda)y)\le \lambda f(x)+(1-\lambda)f(y),\ 0<\lambda<1
\]

\textbf{Interpretation}
If inequality holds for all pairs and all $\lambda$ in interval, function is convex.

\subsection*{N3) Non-dominated sorting rank (mini pattern)}
\textbf{Given}
Front 1: $\{A,B,D,F\}$, Front 2: $\{E,G,H,I\}$, Front 3: $\{C\}$.

\textbf{Find}
Dominance rank assignment.

\textbf{Answer}
\begin{itemize}
\item Rank 1: all members of Front 1.
\item Rank 2: all members of Front 2.
\item Rank 3: all members of Front 3.
\end{itemize}
Lower rank indicates higher preference in Pareto methods.

\section*{Solved Questions for Module 5 (Self-Contained)}

\subsection*{Part A / 3-Mark Questions (Each Answer Has 6 Scoring Points)}

\textbf{Q1. Differentiate linear and nonlinear MOOP.}
\begin{enumerate}
\item Linear MOOP has all objective and constraint functions linear.
\item Nonlinear MOOP has at least one nonlinear objective or constraint.
\item Linear MOOP often has stronger theoretical properties.
\item Nonlinear MOOP is common in practical real-world problems.
\item Nonlinear MOOP solution methods are often computationally harder.
\item Convergence guarantees are generally stronger for linear cases.
\end{enumerate}

\textbf{Q2. Differentiate convex and non-convex MOOP.}
\begin{enumerate}
\item Convex MOOP has convex objective behavior and feasible region.
\item In convex optimization, local minimum is global minimum.
\item Convexity satisfies chord inequality condition.
\item Non-convex MOOP may have multiple local minima.
\item Non-convex search is more difficult and can trap local optima.
\item Most practical MOOPs are nonlinear and often non-convex.
\end{enumerate}

\textbf{Q3. Explain dominance in MOOP.}
\begin{enumerate}
\item Dominance compares two multi-objective solutions.
\item A dominates B if A is no worse in all objectives.
\item A must also be strictly better in at least one objective.
\item If neither dominates the other, both can be non-dominated relative to each other.
\item Dominance is the basis for Pareto ranking.
\item Non-dominated solutions represent trade-off candidates.
\end{enumerate}

\textbf{Q4. What are properties of dominance relation?}
\begin{enumerate}
\item Dominance is not reflexive in strict form.
\item Dominance is not symmetric.
\item Dominance is not antisymmetric in strict relation form.
\item Dominance is transitive.
\item These properties drive ranking in non-dominated sorting.
\item Transitivity enables layered Pareto-front construction.
\end{enumerate}

\textbf{Q5. What is Pareto optimality?}
\begin{enumerate}
\item Pareto-optimal solution is not dominated by any feasible solution.
\item Pareto set is collection of all such non-dominated feasible solutions.
\item Pareto front is image of Pareto set in objective space.
\item Moving along front improves one objective at expense of another.
\item No single universal best point exists for conflicting objectives.
\item Decision maker chooses preferred compromise from Pareto front.
\end{enumerate}

\textbf{Q6. State characteristics of neuro-fuzzy hybrid systems.}
\begin{enumerate}
\item Combines fuzzy reasoning with neural learning.
\item Can process both numerical and linguistic information.
\item Handles imprecise and partial information effectively.
\item Supports self-learning and self-tuning behavior.
\item Uses layered architecture with fuzzy-rule representation.
\item Produces crisp outputs after fuzzy inference and defuzzification.
\end{enumerate}

\subsection*{Part B / 14-Mark Questions (Each Answer Has 14 Scoring Points)}

\textbf{Q1. Explain convex and non-convex MOOP. How to find a non-dominated set?}
\begin{enumerate}
\item MOOP optimizes multiple often-conflicting objectives simultaneously.
\item Convex MOOP obeys convexity inequality over feasible region.
\item In convex problems, local optimum equals global optimum.
\item Non-convex MOOP violates convexity condition in at least part of space.
\item Non-convex problems can contain many local optima.
\item Non-dominated set is built using pairwise dominance tests.
\item For each solution, compare against all others in feasible candidate set.
\item Mark a solution dominated if another is no worse in all objectives and better in at least one.
\item Remove all dominated solutions.
\item Remaining solutions form non-dominated set for that candidate pool.
\item If candidate pool is whole feasible space, this becomes Pareto-optimal set.
\item Mapped curve of this set in objective space is Pareto front.
\item In practice, non-dominated sorting gives front-wise ranking.
\item Conclusion: convexity affects landscape difficulty, while non-dominated filtering identifies trade-off-optimal solutions.
\end{enumerate}

\textbf{Q2. Explain Genetic-Neuro hybrid system in detail with block diagram and advantages.}
\begin{enumerate}
\item Genetic-neuro hybrid combines GA search with ANN learning model.
\item GA maintains a population of candidate ANN parameter sets.
\item Selection chooses promising parents.
\item Crossover recombines parents into offspring.
\item Mutation injects random diversity.
\item Offspring parameters are fed to ANN for evaluation.
\item ANN computes performance/fitness (e.g., inverse error).
\item Fitness values guide next GA generation.
\item Iteration continues until target performance is reached.
\item GA can optimize ANN weights.
\item GA can optimize topology (layers, nodes, interconnections).
\item GA can optimize control parameters like learning rate and momentum.
\item Advantages include global-search support and improved ANN design robustness.
\item Conclusion: hybridization combines GA exploration with ANN modeling power.
\end{enumerate}

\textbf{Q3. What are classifications of neuro-fuzzy hybrid systems? Discuss in detail.}
\begin{enumerate}
\item Neuro-fuzzy systems are hybrids where fuzzy system is trainable via neural principles.
\item Typical layered view: input layer, rule layer, output layer.
\item Crisp inputs enter first layer directly.
\item Fuzzification neurons compute membership grades.
\item Rule layer combines fuzzy evidence according to rule structure.
\item Aggregation combines rule outputs.
\item Defuzzification provides final crisp output.
\item Classification in notes context is based on hybrid architecture family (neuro-fuzzy as one major hybrid type).
\item Hybrid-system type list includes neuro-fuzzy, neuro-genetic, fuzzy-genetic.
\item Neuro-fuzzy key characteristic: local learning causes local model changes.
\item It supports uncertain/vague information processing.
\item It supports self-learning, self-organizing, and self-tuning behavior.
\item Limitations include model-development difficulty and membership-value design challenge.
\item Conclusion: neuro-fuzzy systems provide interpretable fuzzy reasoning with adaptive neural training capability.
\end{enumerate}

\textbf{Q4. Explain Pareto optimality in detail.}
\begin{enumerate}
\item In MOOP, objectives often conflict; one-point optimum is rare.
\item A solution is Pareto-optimal if no feasible solution dominates it.
\item Dominance means no worse in all objectives and better in at least one.
\item Set of Pareto-optimal solutions forms Pareto-optimal set.
\item Image of this set in objective space is Pareto front.
\item Front points are compromise solutions.
\item Improving one objective on front generally worsens another objective.
\item Dominated solutions lie behind or inside front region.
\item Unattainable ideal points lie beyond feasible front.
\item Non-dominated sorting provides practical approximation of front.
\item Rank-1 front corresponds to best non-dominated candidates.
\item Lower-rank fronts are progressively dominated layers.
\item Good MOOP algorithm seeks convergence and diversity on front.
\item Conclusion: Pareto optimality is the central decision framework for conflicting objective optimization.
\end{enumerate}

\textbf{Q5. What are properties of dominance relation?}
\begin{enumerate}
\item Dominance is binary relation between solution pairs.
\item It is not reflexive in strict dominance form.
\item It is not symmetric.
\item It is not antisymmetric in strict comparison interpretation.
\item It is transitive.
\item Transitivity supports front-by-front decomposition.
\item Reflexive failure arises because strict-better condition disallows self-dominance.
\item Symmetry failure is natural in preference ordering.
\item These properties define search pressure in Pareto algorithms.
\item Dominance avoids scalar-weight bias in multi-objective comparison.
\item It works for minimization/maximization with proper sign convention.
\item It underpins non-dominated set extraction.
\item It also supports rank-based evolutionary selection.
\item Conclusion: dominance properties make Pareto optimization mathematically consistent and algorithmically practical.
\end{enumerate}

\section*{Formula Sheet (Module 5)}
\begin{itemize}
\item MOOP: $\min/\max\ f_i(x),\ i=1,\dots,m$
\item Constraints: $g_j(x)\le0$, $h_k(x)=0$
\item Convexity: $f(\lambda x+(1-\lambda)y)\le \lambda f(x)+(1-\lambda)f(y)$
\item Dominance: $x^1\prec x^2$ if $x^1$ no worse in all objectives and better in at least one
\item Pareto set: non-dominated set of whole feasible space
\item Pareto front: mapped boundary of Pareto set in objective space
\end{itemize}

\section*{One-Page Quick Revision Summary}
\textbf{Memory hooks}
\begin{itemize}
\item MOOP gives a set of trade-off optima, not one universal optimum.
\item Dominance is the language of comparison in MOOP.
\item Pareto front = best compromise curve.
\item Genetic-neuro = GA optimizes ANN parameters.
\item Neuro-fuzzy = fuzzy reasoning + neural learning.
\end{itemize}

\textbf{Rapid answer skeleton}
\begin{enumerate}
\item Write definition and condition equation.
\item Draw tiny trade-off/block diagram.
\item List properties/steps.
\item Give one mini example.
\item End with interpretation line.
\end{enumerate}

\examtip{For Pareto questions, draw a tiny 2-objective sketch and mark dominated vs non-dominated points; this is an easy way to secure explanation marks quickly.}

\end{document}
